\section{Discussion and Conclusion}

\subsection{Why This Paper Matters}
The broader contribution of this work is to reframe multi-subject personalized generation as an evaluation problem rather than only a generation problem. If the community cannot reliably measure subject omission, appearance entanglement, and interaction errors, progress in model design will remain difficult to interpret. BindBench and AnchorScore aim to provide the missing measurement infrastructure.

\subsection{Current Limitations}
The proposed benchmark and metric design still involve several limitations. First, to establish a fair and rigorously controlled ablation across various backbone sizes (0.8B to 4B) and fine-tuning paradigms (layer-only vs. LoRA), the current AnchorScore models were trained under a strict compute-constrained budget (600 optimizer steps, observing roughly 9.6K samples). However, BindBench provides up to 60K full silver training pairs. Second, the silver set depends on teacher VLMs, so any systematic teacher bias may propagate into training even after agreement filtering. Third, the diagnostic labels are intentionally coarse and operate at the image level rather than the per-subject level. This is a deliberate trade-off for annotation stability, but it may limit the granularity of downstream analysis. Fourth, the current training set is narrower in generator coverage than the human evaluation set, so strong transfer on unseen generators must be demonstrated empirically rather than assumed. Fifth, the current setup focuses on pairwise evaluation under controlled prompts and reference subjects; broader generalization to arbitrary open-world scenes remains to be studied.

\subsection{Future Work}
Several next steps follow naturally. A critical direction is to explore the \textit{Scaling Laws} of multi-subject evaluators: investigating the absolute performance ceiling of AnchorScore when fully trained on the entire 60K dataset without compute restrictions. Another direction is to study whether the learned evaluator can generalize across unseen generators without retraining. A third is to extend the benchmark toward video or temporally consistent multi-subject generation. A fourth is to investigate more structured diagnostic heads once sufficient evidence is available that the simpler binary scheme has saturated.

\subsection{Conclusion}
We presented the conceptual and methodological foundation for BindBench and AnchorScore, a benchmark-and-metric framework aimed at evaluating multi-subject personalized generation under the failure modes that matter most in practice. The benchmark emphasizes controlled pairwise comparison, a large dual-teacher silver training pipeline, a broader multi-generator human evaluation layer, and image-level diagnostic supervision. The metric couples pairwise ranking with interpretable error prediction in order to preserve score separability as subject count increases. With the experimental section completed, this framework is positioned to test the central claim of the paper: that existing metrics collapse in the high-subject regime, while a reference-conditioned evaluator can remain discriminative and human-aligned.
